\documentclass{article}
\usepackage{decomp}
\begin{document}
\sourcefile{yeatts2026\_supplement.pdf}
\decompnote{``Supplemental Methods and Results'' of the bioRxiv preprint (media-1.pdf), 5 pages, no bioRxiv stamp. Left out on purpose: the margin line numbers (1--54) and the page numbers. Each supplemental figure is placed before its caption, where it is printed. The text was extracted from the PDF's text layer by work/build\_main.py and proofread against the page images.}

\pdfpage{1}
\subsection*{\underline{Supplemental Methods and Results}}

Further detail on the limitations of the previous imaging setup, background subtraction methodology, and optical components are provided below in the supplemental figures 1,2, and 4, respectively. Additionally, the 3D slide deck of subject V is visualized similar to Figure 6 for subject M in supplemental figure 3. The customization of the VF-800 compresstome is elaborated and alternative gel fixation methods explored to further section integrity are detailed below.

\subsubsection*{Compresstome Customization}

Once we received the compresstome, we worked with engineers from the manufacturing company (Precisionary Instruments) to improve and test blade holder designs addressing corrosion from the PBS on the oscillating mechanism and section movement during cutting. A splash guard was added to the oscillating mechanism, and the front of the blade holder was curved and windowed. Additionally, a cooling system was implemented, with a Lytron Kodiak series recirculating chiller to regulate the sample temperature. The chiller was set to 8 °C while sample bath temperature was fluctuating around 10°C. Slice quality and sample integrity improved with these changes, which were implemented between Subject M and Subject V. Sections were collected intact from the imaging bath and placed into long-term storage as visualized in figure 2d.

\subsubsection*{Alternative Gelling Investigation}

To maintain compatibility with downstream tissue clearing, we investigated the feasibility of hydrogel embedding using the HMS formulation without paraformaldehyde (PFA). Standard SMART protocols include PFA in the hydrogel monomer solution, where it crosslinks tissue biomolecules and anchors the polymer network within the sample\textsuperscript{3}. However, concerns about over-fixation and its potential impact on clearing efficiency motivated efforts to reduce or eliminate PFA during gelation.

We first evaluated bulk gelation in the absence of tissue using the HMS formulation without PFA. To optimize gel formation, polymerization conditions were systematically varied, including temperature and duration. Although gelation could be achieved under extended conditions, the resulting hydrogels were mechanically weak and inconsistent. When these conditions were applied to test tissue, the hydrogel failed to adhere and did not form a stable, tissue-integrated \pdfpage{2} matrix. While acrylamide polymerization likely occurred, the absence of PFA prevented effective covalent coupling between the forming hydrogel and tissue components.

These results are consistent with the underlying chemistry of CLARITY-type methods. PFA forms methylene bridge crosslinks between proteins and facilitates interactions between tissue biomolecules and acrylamide monomers, enabling formation of a stable hydrogel–tissue hybrid. In its absence, the polymer network remains unanchored, resulting in poor mechanical integrity and loss of structural support during handling\textsuperscript{39}. Modifying polymerization conditions alone was insufficient to compensate for the lack of chemical crosslinking.

Together, these findings indicate that inclusion of PFA, or an alternative crosslinking strategy, is required for effective hydrogel embedding in large brain samples, even when downstream applications impose constraints on fixation.

\begin{center}
\includegraphics[width=0.95\linewidth,height=0.75\textheight,keepaspectratio]{figures/sfig1.png}
\end{center}
\noindent \textbf{Supplemental Figure 1: Vibratome Limitations. a.} Fixed optics and a moving sample. (1) referring to the imaging position with a maximum translation of 20 mm in the x and y (red and yellow arrows). (2) slicing position to engage with the vibratome blade \textbf{b.} Motor controlled aluminum stage. Agar gel backing to prevent blade damage from hitting the stage during slicing. Sample of a blocked section of macaque brain around the thalamus. \textbf{c.} The maximum blade range is the total cutting distance of the vibratome. The maximum sample height is determined by the distance from the stage surface, where the sample is glued, to the blade. \textbf{d.} Fixed sample and moving optics \textbf{e.} Photo of Subject M between data acquisition. \textbf{f.} The \pdfpage{3} maximum sample diameter is less than the opening of the specimen tube to ensure the sample does not touch the specimen tube walls. The maximum sample height is the distance from the cutting surface (top of the specimen tube) and the plunger surface in the lowest position.

\alttext{Six panels (a--f), one raster image (1908 × 1137 px). Panel a: a 3D schematic of the previous vibratome set-up: a dark plate with a lens (the fixed optics) above a transparent bath; a red beam comes down onto a brown tissue block in yellow gel at position ``(1)'', which carries three double arrows (red and yellow horizontal for x and y, blue vertical); a second block at ``(2)'', to the right, sits against the dark blade holder, and a bent arrow shows the shuttle path between the two positions. Panel b: a rendered view of the vibratome stage: a grey L-shaped block labelled ``Stage'' (white arrowhead), a light-blue layer labelled ``Agar gel'' (arrowhead) behind a brown block labelled ``Sample'' showing a cut face with white-matter pattern. Panel c: a side-view schematic of the vibratome bath: the blade holder (grey) above the water, a horizontal red double arrow with dashed lines labelled ``Maximum Blade Range'' spanning the blade's travel, and a blue vertical double arrow labelled ``Maximum Sample Height'' from the stage surface to the blade, next to a small tissue block. Panel d: a 3D schematic of the new set-up: a square dark-blue bath with the round top of the specimen tube (a light disc) holding a brain section; the scan head (dark plate with lens) and its vertical post are drawn in several positions (solid, and translucent pink and yellow copies) with red, yellow and blue double arrows, i.e. the optics move in x, y and z while the sample stays. Panel e: a photograph of a coronal macaque brain section (Subject M, pale cream tissue with the white matter slightly lighter, the ventricles and basal ganglia visible) embedded in a grey-blue agar block; labels ``Agar gel'' (top left) and ``Sample'' (on the tissue); a white scale bar ``10 mm'' at the bottom right. Panel f: a cut-away schematic of the specimen tube (grey cylinder) with a brain in sagittal view standing on the white plunger disc; a dark bar across the tube's top represents the blade; ``Maximum Sample Diameter'' (red, with a red double arrow across the opening and dashed verticals) and ``Max Sample Height'' (blue vertical double arrow from the cutting surface down to the plunger).}

\begin{center}
\includegraphics[width=0.95\linewidth,height=0.75\textheight,keepaspectratio]{figures/sfig2.png}
\end{center}
\noindent \textbf{Supplemental Figure 2: Background Subtraction. a.} Raw spectral data amplitude for a given camera pixel of one A-line, blue representing channel 2 and orange channel 1. \textbf{b.} The amplitude of the complex depth profile (the pre-processed A-line) against depth of the same A-line. \textbf{c.} Depth profiles in dB scale without and with background subtraction.

\alttext{One raster image (768 × 652 px, only 121 ppi) with three panels. Panels a and b are each two stacked line plots in thin grey axes boxes (left y axis, bottom x axis, short unlabelled tick marks, no numbers on any axis), two series per plot: blue (channel 2) and orange (channel 1). Panel a (left): y label ``Intensity'' shared by both boxes, x label ``Pixel number''; 10 evenly spaced x ticks (every 0.098 of the box width). Top box: the two raw spectra, both starting near 0 at the left, rising steeply, blue earlier than orange; blue peaks at about 0.19 of the width at 0.78 of the box height (digitized), orange at 0.20 at 0.84 with a ripple of about ±0.05; both dip to about 0.52 near 0.37, rise to a broad second hump of about 0.75 around 0.60--0.70, and fall to about 0.05 at the right edge; orange is slightly above blue on the right slope. Bottom box: the same spectra after background subtraction, an interferogram oscillating about a baseline at 0.45 of the height; the orange trace swings widely (envelope from about 0.0 to 1.0 near x = 0.24, and beats with large amplitude near x = 0.2--0.3, 0.5 and 0.8), the blue trace swings little (envelope about 0.3--0.6). Panel b (right): y label ``Intensity (dB)'', x label ``Depth''; 5 evenly spaced x ticks. Top box (without background subtraction): blue starts high at the left (0.58 of the height) and decays smoothly to about 0.36 by x = 0.2 (the background), orange starts at 0.54 and runs noisily around 0.3 with dips; at x = 0.25 both jump (the tissue surface) to peaks of 0.70 (blue) and 0.79 (orange) at x = 0.27, then decay noisily (orange above blue by about 0.1) to about 0.2--0.3; a sharp narrow peak at x = 0.86 reaches 0.68 (blue) and 0.76 (orange); both end near 0.2 at the right. Bottom box (with background subtraction): flat noisy baselines at the left (blue 0.37, orange 0.45), the same surface rise at 0.25 to 0.75 (blue) and 0.84 (orange), the same slow decay, the narrow peak at 0.86 (0.56 blue, 0.81 orange) and ends near 0.4. Panel c: two wide grey-scale B-scan images (each 184 × ~713 px, lateral position across, depth downwards), each with a vertical grey bar at its right running from white (top) to black (bottom) without numbers. Top image (without background subtraction): the upper quarter is filled with bright horizontal banding (a hazy, laterally uniform background that is brightest at the very top and fades downwards), then the tissue surface as an undulating bright line at about a quarter of the depth, bright speckled patches (tissue) below it fading with depth, a thin bright horizontal line across the whole width at about 85\% of the depth and a dark band beneath it. Bottom image (with background subtraction): the same scene with the upper band now black, so the surface line and the patches stand out; the thin bright line near the bottom remains.}
\figuredata{a, top box}{figures/sfig2atop.csv}{figures/sfig2a.py}{tier 3, digitized from the 121-ppi embedded raster (colour masks, per-column envelope); no axis numbers, coordinates are axis-box fractions}
\figuredata{a, bottom box}{figures/sfig2abottom.csv}{figures/sfig2a.py}{tier 3, digitized as above}
\figuredata{b, top box}{figures/sfig2btop.csv}{figures/sfig2b.py}{tier 3, digitized as above}
\figuredata{b, bottom box}{figures/sfig2bbottom.csv}{figures/sfig2b.py}{tier 3, digitized as above}
\figuredata{c, top image}{figures/sfig2ctop.csv}{figures/sfig2c.py}{tier 3, grey-level grid copied from the embedded raster at native resolution; the grey bar has no numbers, so uncalibrated}
\figuredata{c, bottom image}{figures/sfig2cbottom.csv}{figures/sfig2c.py}{tier 3, grey-level grid as above, uncalibrated}

\pdfpage{4}
\begin{center}
\includegraphics[width=0.95\linewidth,height=0.75\textheight,keepaspectratio]{figures/sfig3.png}
\end{center}
\noindent \textbf{Supplemental Fig. 3: Alignment across sections produces 3D PS-OCT volumes in} \textbf{monkey V.} After imaging in the sagittal plane, axial and coronal views show good between-slice alignment.

\alttext{Two dark grey-scale raster images (1481 × 1308 px and 1481 × 653 px) on black, rather dim. Panel a: three reconstructed views of the aligned sagittal cross-polarization stack of monkey V, side by side: left a midline sagittal section (the imaging plane; corpus callosum arching over the thalamus, the cerebellar arbor vitae bright at the lower left), middle an axial (horizontal) section with the two hemispheres and the basal ganglia, right a coronal section through the basal ganglia; a dark horizontal strip runs along the top of the panel. Panel b: two montages: the top one is 5 rows × about 24 small axial sections that progress from tiny dorsal caps (top row) to full hemispheres (rows 3--4) and back to smaller ventral shapes; the bottom one is 9 rows × about 24 coronal sections from small frontal shapes to full sections and down to tiny occipital/cerebellar pieces. Fine vertical banding from the slice boundaries is visible in all views. No scale bars or labels other than the panel letters.}

\pdfpage{5}
\begin{center}
\includegraphics[width=0.95\linewidth,height=0.75\textheight,keepaspectratio]{figures/sfig4.png}
\end{center}
\noindent \textbf{Supplemental Figure 4: Optical system.} Detailed optical components of the interferometer. SLD: superluminescent diode, PM Coupler: polarization maintaining coupler, PMF: polarization maintaining fiber, QWP: quarter wave plate.

\alttext{A flat-colour optical schematic (one raster image, 1911 × 1771 px) above a legend separated by a thick black line. Fibre paths are thin blue lines (PMF), free-space beams are thick red bars (pink where they are diverted to a calibration path). Top left, ``Source Arm'': the SLD (dark square with a white star-burst) → collimator (light-blue ellipse) → polarizer (tall light-blue plate) → collimator → fibre to the ``PM Coupler'' (dark rounded rectangle in the centre). Top right, ``Reference Arm'': fibre → collimator → beam sampler (a thin tilted plate) that reflects a pink beam up-left through a QWP (small grey plate) and a lens (ellipse) onto a hatched mirror (the calibration path), while the main red beam continues through a polarizer (tall plate) and a lens to a hatched glass wedge at the right. Lower right, ``Sample Arm'': fibre → collimator → beam sampler (again diverting a pink beam up-left through a QWP and a lens to a mirror) → QWP → two crossed galvo mirrors (dark slanted bars) that turn the beam downward → lens (scan lens) → the sample, a brain lying in a tray of liquid. Bottom left, ``Spectrometer'': fibre → collimator → grating (a tilted grey slab) → Wollaston prism (a square split diagonally into two blues) → lens → camera (dark box). Legend, with the symbol drawn before each name: ``SLD'' (dark square with star-burst), ``Glass Wedge'' (hatched bar), ``Beam Sampler'' (thin tilted plate), ``Galvo Mirror'' (two crossed dark bars), ``Grating'' (grey bar), ``Collimator'' (light-blue ellipse), ``Lens'' (thin white ellipse), ``Mirror'' (hatched dark bar), ``Wollaston Prism'' (two-tone blue square), ``PM Coupler'' (dark rounded rectangle), ``Polarizer'' (tall light-blue plate), ``QWP'' (small grey plate), ``PMF'' (blue line), ``Camera'' (dark box).}

\end{document}
